\section{Các nghiên cứu liên quan}

Những năm gần đây, bài toán cải thiện độ chính xác ước tính camera pose của thuật toán SfM truyền thống được quan tâm rộng rãi trong cộng đồng nghiên cứu. Với sự nổi lên của các thuật toán tái tạo 3D mạnh mẽ như NeRF, Gaussian Splatting \cite{kerbl2023gaussian} thì việc ước tính camera pose chính xác sẽ tạo ra một đầu vào chất lượng, giảm thiểu sai số cho quá trình tái tạo, tạo ra các mô hình mây điểm 3D chi tiết, chính xác cao. 

Trong các hướng tiếp cận giải quyết bài toán này, việc sử dụng các đặc trưng điểm học sâu như SuperPoint \cite{detone2018superpoint}, D2-Net \cite{dusmanu2019d2}, ALIKED đang trở thành xu hướng chủ đạo và đã được chứng minh là vượt trội hơn so với các phương pháp truyền thống sử dụng bộ trích xuất đặc trưng SIFT trong nhiều tập dữ liệu nổi tiếng khác nhau. Cụ thể, trong bài báo \cite{geometryaware2024}, các thử nghiệm đã chỉ ra rằng các đặc trưng học sâu này giúp tăng số lượng ảnh được tái thiết lập và giảm sai số chiếu lại trung bình, cải thiện độ chính xác ước lượng camera pose đặc biệt trong các điều kiện ánh sáng thay đổi hoặc cảnh có ít kết cấu. 

Bên cạnh đó, các tiếp cận sử dụng các đặc trưng hình học như đường thẳng, góc vào quy trình SfM đã được chứng minh là có tiềm năng và cải thiện đáng kể độ chính xác và độ ổn định của tái dựng 3D, đặc biệt trong các môi trường có ít kết cấu, đơn giản hoặc cấu trúc lặp lại. Tuy nhiên, \cite{Liu_2023_CVPR} chỉ ra rằng, mặc cho việc mang nhiều đặc trưng tiêu biểu về cấu trúc của khung cảnh, thì kết quả thực nghiệm của các phương pháp tiếp cận dựa trên đặc trưng đường thẳng vẫn có kết quả kém hơn so với các phương pháp dựa trên điểm. Các nghiên cứu gần đây đã đề xuất các phương pháp kết hợp đặc trưng điểm và đường thẳng để tăng cường ràng buộc hình học trong quá trình tái dựng. Ví dụ, phương pháp SfM kết hợp đặc trưng lai ghép \cite{liu2024robust} đã được phát triển để tận dụng cả điểm và đường thẳng, giúp cải thiện độ chính xác của mô hình 3D và ước tính camera pose, đặc biệt trong các cảnh có cấu trúc mạnh như kiến trúc đô thị. 

Mặc dù, cách tiếp cận dựa trên kết hợp điểm và đường thẳng có rất nhiều tiềm năng nhưng nhìn chung các phương pháp dựa trên đường thẳng này thường không phù hợp đối với các bộ dữ liệu bay chụp cột ăng-ten, trạm viễn thông bởi các lí do sau đây. Một là, địa hình lắp đặt cột rất đa dạng bao gồm: trên mái nhà, cánh đồng, rừng núi, biển,... Trong các địa hình này, chỉ có khu vực đô thị là có giàu các đường nét, góc cạnh từ các ngôi nhà, các địa hình còn lại thường có rất ít các đặc trưng như vậy. Do đó việc áp dụng phát hiện các đặc trưng đường thẳng là không phù hợp để áp dụng vào thực tế. Hai là, cấu tạo của cột ăng-ten thường là các hình trụ thẳng đứng với bề mặt trơn nhẵn, đồng màu. Cho nên việc phát hiện các đường thẳng trên cột là kém hiệu quả và dễ gây nhiễu khi quan sát các ảnh 2D với các góc nhìn khác nhau có thể tạo ra các đường thẳng đặc trưng ảo trên ảnh.

Trong phạm vi nghiên cứu của luận văn, luận văn tập trung vào nghiên cứu các phương pháp cải tiến dựa trên việc sử dụng các điểm đặc trưng học sâu - hướng tiếp cận được cho là tốt nhất hiện tại trong lĩnh vực tái tạo 3D \cite{Liu_2023_CVPR}. Xuất phát từ ý tưởng sử dụng các điểm đặc trưng dựa trên học sâu, Hybrid-SfM được đề xuất với giải pháp kết hợp (chứ không thay thế) điểm mạnh của đặc trưng truyền thống và điểm mạnh của đặc trưng học sâu để phù hợp, hiệu quả cho bài toán thực tế đối với dữ liệu bay chụp cột ăng-ten, trạm viễn thông. 

% Tái tạo 3D là một công nghệ quan trọng với các ứng dụng thực tiễn trong viễn thông, nông nghiệp, và điện lực, nhờ khả năng tạo mô hình không gian chính xác từ dữ liệu hình ảnh hoặc cảm biến. Dưới đây là các nghiên cứu tiêu biểu về tái tạo 3D, với mô tả chi tiết về cách xây dựng mô hình 3D trong từng phương pháp.

% Trong lĩnh vực viễn thông, \cite{pastucha2023} đề xuất tái tạo hành lang dây cáp bằng hình ảnh từ máy bay không người lái (UAV). Mô hình 3D được xây dựng bằng thuật toán SfM, sử dụng khoảng 100 ảnh RGB chụp từ nhiều góc độ với độ chồng lấn 80\%. Các ảnh được xử lý để tạo đám mây điểm 3D thông qua khớp đặc trưng (SIFT), cho phép phát hiện chướng ngại vật vi phạm hành lang với sai số dưới 10 cm.

% Trong nông nghiệp, \cite{zhu2024} giới thiệu tập dữ liệu Crops3D để tái tạo 3D cây trồng như cà chua và bắp cải. Mô hình 3D được tạo bằng hai cách: (1) SfM từ khoảng 100 ảnh RGB chụp vòng quanh cây, sử dụng phần mềm như COLMAP để khớp đặc trưng và tạo đám mây điểm; (2) quét laser bằng máy ánh sáng cấu trúc, thu thập dữ liệu điểm 3D trực tiếp với mật độ cao. Kết quả hỗ trợ phân tích hình thái cây và phân đoạn cây trồng với độ chính xác trên 90\%.

% Trong điện lực, \cite{darwish2023} phát triển PL-NeRF, một phương pháp dựa trên NeRF, để tái tạo đường dây điện từ hình ảnh chụp bởi rô-bốt kiểm tra (FPLIR). Mô hình 3D được xây dựng bằng cách tối ưu hóa trường bức xạ thần kinh từ chuỗi ảnh liên tục, sử dụng mã hóa vị trí tích hợp và mã hóa băm để xử lý cấu trúc mỏng của dây điện.

% Bên cạnh đó, \cite{xiao2024} tập trung vào tái tạo 3D cotton bolls trong nông nghiệp bằng UAV. Mô hình 3D được tạo bằng SfM, sử dụng hàng trăm ảnh RGB chụp từ độ cao 10--20 m. Các ảnh được xử lý bằng phần mềm Agisoft Metashape để khớp đặc trưng, tạo đám mây điểm, và xây dựng lưới 3D, hỗ trợ phân tích đặc điểm cây trồng với độ phân giải dưới 1 cm.

% Các phương pháp trên cho thấy sự đa dạng trong tái tạo 3D dựa trên hình ảnh. Mỗi phương pháp có ưu điểm riêng, phù hợp với yêu cầu cụ thể của viễn thông, nông nghiệp, và điện lực, thúc đẩy quản lý hạ tầng và tối ưu hóa sản xuất. Phần lớn các ứng dụng thực tế hiện tại sử dụng quy trình xây dựng đám mây điểm 3D truyền thống là kết hợp SfM (cho xây mây điểm thưa) và MVS (cho xây dựng mây điểm chi tiết). Trong Chương 3, đối với bài toán tái tạo mây điểm 3D chi tiết ứng dụng thực tế cho cột ăng-ten của trạm viễn thông, luận văn sẽ đề xuất kỹ thuật cải thiện độ chính xác khi tính toán camera pose trong quy trình SfM và khảo sát, đánh giá mô hình mây điểm chi tiết trên 2 phương pháp truyền thống MVS và phương pháp hiện đại NeRF.
